\documentclass[12pt,a4paper]{article}
\usepackage[utf8]{inputenc}
\usepackage[T1]{fontenc}
\usepackage[italian]{babel}
\usepackage{geometry}
\usepackage{longtable}
\usepackage{array}
\usepackage{booktabs}
\usepackage{hyperref}
\usepackage{xcolor}
\usepackage{listings}
\usepackage{enumitem}
\usepackage{tikz}

\geometry{margin=2cm}

\lstdefinestyle{javastyle}{
  basicstyle=\ttfamily\small,
  keywordstyle=\color{blue},
  commentstyle=\color{gray},
  stringstyle=\color{red},
  breaklines=true,
  showstringspaces=false,
  frame=single
}

\title{Spiegazione Approfondita delle Funzioni\\Progetto Concessionario Auto \& Moto\\Architettura BCE + DAO}
\author{N86005495 \and N86005712}
\date{\today}

\begin{document}

\maketitle
\newpage
\tableofcontents
\newpage

\section{Introduzione e Struttura del Documento}

Questo documento spiega il progetto seguendo una struttura per \textbf{funzioni specifiche}, non per codice copiato.
L'obiettivo: capire \textbf{cosa fa} ogni funzione, \textbf{come è utilizzata}, e \textbf{perché esiste}.

Ogni funzione è descritta con:
\begin{enumerate}
  \item \textbf{Cosa fa}: responsabilità concreta della funzione.
  \item \textbf{Come viene usata}: in quale contesto, da quale altro codice.
  \item \textbf{Perché esiste}: problema che risolve, valore che aggiunge all'architettura.
  \item \textbf{File}: dove si trova nel progetto.
\end{enumerate}

Le funzioni sono organizzate per \textbf{livelli architetturali} (Entry Point $\to$ GUI $\to$ Controller $\to$ DAO $\to$ Entity $\to$ Database), non per file singoli.
Se una funzione è usata in modi diversi, è spiegata per ogni contesto senza ripetizione dei concetti base.

\section{Architettura complessiva: il pattern BCE + DAO}

Il progetto separa il codice in cinque livelli rigidamente separati:

\begin{enumerate}
  \item \textbf{Boundary (B):} GUI, widget Swing. Responsabilità: presentare informazioni all'utente, catturare input, delegare al controller.
    Non contiene logica di business, validazione complessa, o accesso diretto a DB.
    Se il controller cambia, la GUI non deve cambiare (teoricamente).
  
  \item \textbf{Control (C):} Orchestratore della logica. Responsabilità: ricevere richieste dalla GUI, validare input, coordinare i DAO, applicare regole di business, tradurre eccezioni tecniche in semantiche applicative.
    Il controller è l'unico punto di contatto tra GUI e DAO: la GUI non chiama direttamente il DAO.
  
  \item \textbf{Entity (E):} Modelli di dominio immutabili (Cliente, Veicolo, Dipendente).
    Rappresentano concetti del mondo reale. Campi immutabili (final), niente getter/setter, niente logica complessa.
    Sono le strutture dati che circolano tra controller e GUI.
  
  \item \textbf{DAO:} Data Access Object. Layer di accesso ai dati che incapsula SQL e JDBC.
    Fornisce un'interfaccia semplice (inserisci, leggi, aggiorna, cancella) al controller senza esporre i dettagli SQL.
    Se il DB cambia (PostgreSQL $\to$ MongoDB), scrivi una nuova implementazione DAO, il controller non sa la differenza.
  
  \item \textbf{Database:} PostgreSQL con vincoli (constraint), funzioni (stored procedure), trigger, e transazioni.
    È la fonte unica di verità: lo stato del DB è lo stato dell'applicazione.
    Applica integrità a livello database: vincoli check, foreign key, trigger per logica che non può essere fidata al livello applicativo.
\end{enumerate}

\textbf{Flusso verticale di un'operazione (esempio: utente clicca ``Salva Cliente''):}

\begin{center}
\begin{tabular}{|l|l|}
\hline
Livello & Cosa accade \\
\hline
Boundary (GUI) & Utente digita dati, clicca bottone. GUI legge form e chiama Controller. \\
\hline
Control (Controller) & Riceve dati dal pannello. Valida (email formato, telefono 10 cifre). \\
& Controlla duplicati (email già registrata?). Se errore $\to$ lancia exception. \\
& Se ok $\to$ chiama DAO. \\
\hline
DAO & Riceve dati validati. Genera SQL INSERT. Apre connessione. Esegue. \\
& Se SQL fallisce (constraint violato) $\to$ lancia DatabaseException. \\
& Se ok $\to$ ritorna success al controller. \\
\hline
Database & Riceve INSERT. Controlla constraint (email UNIQUE, format check). \\
& Se ok $\to$ inserisce. Se fallisce $\to$ exception SQL. \\
\hline
Ritorno & DAO ritorna risultato (success/exception) al Controller. \\
& Controller cattura exception, la traduce. Ritorna al pannello GUI. \\
& GUI mostra dialog: success o messaggio di errore. \\
\hline
\end{tabular}
\end{center}

\textbf{Perché questa separazione è importante:}

\begin{itemize}
  \item \textbf{Testabilità:} Scrivi test che chiama il Controller senza GUI. MockaI il DAO. Niente browser, niente click.
  \item \textbf{Manutenibilità:} Una query SQL cambia? Modifichi solo ClienteDAOImpl. Il Controller e la GUI non sanno.
  \item \textbf{Riusabilità:} Vuoi una API REST che aggiunge clienti? Chiama lo stesso Controller, bypassa la GUI.
  \item \textbf{Scalabilità:} Vuoi cambiare DB? Implementi ClienteDAOMongo, plugghi nel Controller, niente cambia in GUI o Controller.
  \item \textbf{Debugging:} Un bug in una operazione? Sai esattamente dove cercare: GUI $\to$ Controller $\to$ DAO $\to$ DB. Non è caos.
\end{itemize}



\section{Entry Point e Bootstrap}

\subsection{Main.main() -- File: Main.java}

\textbf{Cosa fa:}
Punto di ingresso assoluto del programma. Riceve il controllo dalla JVM e lo delega a MainApp.

\textbf{Come viene usato:}
Quando l'utente avvia l'applicazione (doppio click su .jar o esecuzione da terminale), la JVM cerca il metodo \texttt{public static void main(String[] args)} e lo esegue da qui.

\textbf{Perché esiste:}
Fornisce un'indirezione: se in futuro devi aggiungere logica di bootstrap (parsing argomenti, config da file, check licenza), lo fai qui senza toccare MainApp.

\subsection{MainApp.main() -- File: MainApp.java}

\textbf{Cosa fa:}
Prepara l'ambiente Swing e lancia la finestra principale (\texttt{ConcessionarioFrame}).

Due responsabilità principali:
\begin{enumerate}
  \item Esegue il setup UI sul thread EDT (Event Dispatch Thread).
  \item Imposta il look-and-feel nativo (Windows, macOS, GTK).
\end{enumerate}

\textbf{Come viene usato:}
Chiamata da \texttt{Main.main()}. Assicura che la finestra Swing emerga sul thread corretto.

\textbf{Perché esiste:}
Swing è thread-sensitive: tutti gli aggiornamenti UI devono avvenire su un thread speciale (EDT).
Se crei componenti Swing dal thread principale, otterrai race condition e crash UI.
\texttt{SwingUtilities.invokeLater(() -> \{ ... \})} mette il blocco di codice in coda all'EDT.
Quando EDT è libero, lo esegue in un contesto thread-safe.

Il \texttt{UIManager.setLookAndFeel()} adatta la GUI al SO (Windows button diverse da macOS, colori GTK su Linux).
Senza questo, l'app ha un look ``generico Java'' che sembra fuori posto su qualsiasi SO.

\section{Livello Boundary: La GUI con Swing}

\subsection{ConcessionarioFrame() -- Constructor. File: ConcessionarioFrame.java}

\textbf{Cosa fa:}
Costruisce la finestra principale. Divide lo spazio in due zone: header statico (titolo + logout) e area dinamica con CardLayout.

\textbf{Come viene usato:}
Istanziata da MainApp. Durante costruzione, prepara il layout, crea il CardLayout, e aggiunge il login panel.

\textbf{Perché esiste:}
È il contenitore di tutta l'UI. All'avvio, mostra login. Dopo autenticazione, sostituisce il content con l'area DIPENDENTE o CLIENTE.
Mantenendo una sola JFrame, non crei finestre spurie; tutto è dentro una finestra unica.

\textbf{CardLayout: il meccanismo di switch:}
Il CardLayout è un layout manager che tiene molteplici pannelli (``card'') sovapposti, mostrando uno alla volta.
Quando utente fa login, chiami \texttt{cardLayout.show(cardsPanel, "CLIENTE\_CARD")}, e il pannello CLIENTE diventa visibile.
Vantaggio: non ricostruisci l'UI; mantieni lo stato.

\subsection{ConcessionarioFrame.createHeader() -- File: ConcessionarioFrame.java}

\textbf{Cosa fa:}
Crea il pannello header (zona NORTH): titolo a sinistra, info sessione + logout a destra.

\textbf{Come viene usato:}
Chiamata da \texttt{ConcessionarioFrame()}, aggiunta al frame con \texttt{add(..., BorderLayout.NORTH)}.
L'header è statico: non cambia durante l'uso dell'app.

\textbf{Perché esiste:}
Separa visualmente il contenuto principale (CENTER) da info di sessione.
Header mostra always chi è autenticato, consente quick logout.

\subsection{ConcessionarioFrame.createLoginPanel() -- File: ConcessionarioFrame.java}

\textbf{Cosa fa:}
Costruisce pannello login con due sezioni: Cliente (email/password) e Dipendente (id/password).
Bottoni separati: ``Login Cliente'', ``Login Dipendente''.

\textbf{Come viene usato:}
Aggiunta al CardLayout come prima card (LOGIN\_CARD).
Utente digita credenziali, clicca bottone $\to$ trigger action listener $\to$ \texttt{handleClienteLogin()} o \texttt{handleDipendenteLogin()}.

\textbf{Perché esiste:}
Punto di ingresso: due percorsi autenticazione (cliente vs. dipendente) con GUI diversa.

\subsection{ConcessionarioFrame.handleClienteLogin() -- File: ConcessionarioFrame.java}

\textbf{Cosa fa:}
Estrae email/password da form, chiama \texttt{Controller.autenticaCliente()}.
Se autenticazione succede, aggiorna header (sessionLabel, abilita logout), costruisce area cliente, mostra card CLIENTE.
Se errore, mostra dialog di errore con messaggio amichevole.

\textbf{Come viene usato:}
Listener del bottone ``Login Cliente''.

\textbf{Perché esiste:}
Gestisce il flusso login: input $\to$ controller $\to$ result $\to$ update header e card.

\subsection{ConcessionarioFrame.buildClienteArea() -- File: ConcessionarioFrame.java}

\textbf{Cosa fa:}
Crea il JTabbedPane con 4 tab per il cliente:
\begin{enumerate}
  \item Profilo: edit email/telefono.
  \item Catalogo: vedi veicoli disponibili (read-only).
  \item Prenota Test Drive: form per richiedere test drive.
  \item Le Mie Operazioni: storico vendite e test drive.
\end{enumerate}

\textbf{Come viene usato:}
Chiamata da \texttt{handleClienteLogin()} dopo autenticazione riuscita.
Pannello ritornato è aggiunto al CardLayout e reso visibile.

\textbf{Perché esiste:}
Organizza le funzionalità cliente in tab logici. Ogni tab è un pannello specializzato (ProfiloClientePanel, CatalogoClientePanel, etc.).

\subsection{ConcessionarioFrame.handleLogout() -- File: ConcessionarioFrame.java}

\textbf{Cosa fa:}
Ripristina lo stato iniziale: mostra LOGIN card, resetta sessionLabel a ``Non autenticato'', disabilita bottone logout.

\textbf{Come viene usato:}
Listener del bottone Logout nell'header.

\textbf{Perché esiste:}
Chiude la sessione corrente senza chiudere l'applicazione.

\section{Livello Boundary: Pannelli Specifici}

\subsection{ClientiPanel -- File: ClientiPanel.java (Dipendente)}

\textbf{Cosa fa:}
Pannello CRUD per dipendenti. Tabella con lista clienti a sinistra, form (nome, email, telefono, tipo) a destra.
Bottoni: Aggiungi, Modifica, Elimina, Aggiorna.

\textbf{Come viene usato:}
Istanziata e aggiunta a JTabbedPane quando dipendente autentica, in tab ``Clienti''.

\textbf{Metodi interni principali:}

\begin{itemize}
  \item \texttt{loadClienti()}: chiama \texttt{Controller.getClienti()}, riceve lista, popola tabella.
  \item \texttt{onRowSelected()}: utente clicca riga in tabella $\to$ estrae dati e riempie form per modifica/view.
  \item \texttt{onAggiungi()}: legge form, valida, chiama \texttt{Controller.inserisciCliente()}, ricarica tabella.
  \item \texttt{onModifica()}: estrae ID da riga selezionata, chiama \texttt{Controller.aggiornaCliente()}.
  \item \texttt{onElimina()}: richiede conferma, chiama \texttt{Controller.eliminaCliente()}.
\end{itemize}

\textbf{Perché esiste:}
Consente al dipendente di gestire anagrafe clienti senza accesso diretto a DB o SQL.

\subsection{VeicoliPanel -- File: VeicoliPanel.java (Dipendente)}

\textbf{Cosa fa:}
CRUD veicoli. Tabella (targa, marca, modello, tipo, prezzo, stato), form per inserimento.

Logica speciale: se la marca non esiste in DB, chiede al dipendente di crearla.

\textbf{Metodi interni:}
\begin{itemize}
  \item \texttt{onAggiungi()}: estrae dati, controlla se marca esiste, se no popup ``Creare marca?''.
    Se sì, chiede nazione di origine e crea marca, poi veicolo.
\end{itemize}

\textbf{Perché esiste:}
Gestione dell'inventario: aggiungi auto, moto, e loro stati (disponibile, venduto).

\subsection{VenditaPanel -- File: VenditaPanel.java (Dipendente)}

\textbf{Cosa fa:}
Form per registrare una vendita: combo clienti, combo veicoli disponibili, campo prezzo finale.
Quando selezioni un veicolo, il prezzo si popola automaticamente dal DB.
Bottone ``Effettua Vendita'' chiama \texttt{Controller.effettuaVendita()}.

\textbf{Perché esiste:}
Registra transazioni di vendita. Il controller controlla che veicolo sia disponibile, dipendente è autenticato, etc.

\subsection{TestDrivePanel -- File: TestDrivePanel.java (Dipendente)}

\textbf{Cosa fa:}
Form per prenotare test drive per un cliente.
Combo cliente, combo veicoli, campo data (future date picker).

\textbf{Come viene usato:}
Tab ``Test Drive'' nell'area dipendente.

\textbf{Perché esiste:}
Consente al dipendente di creare prenotazioni test drive per clienti.

\subsection{ProfiloClientePanel -- File: ProfiloClientePanel.java (Cliente)}

\textbf{Cosa fa:}
Mostra profilo cliente (nome, email, telefono).
Consente edizione email e telefono con bottone ``Salva''.

\textbf{Come viene usato:}
Tab ``Profilo'' nell'area cliente.

\textbf{Perché esiste:}
Consente al cliente di aggiornare i propri contatti.

\subsection{CatalogoClientePanel -- File: CatalogoClientePanel.java (Cliente)}

\textbf{Cosa fa:}
Mostra tabella read-only di veicoli disponibili (targa, marca, modello, tipo, prezzo, stato).

\textbf{Come viene usato:}
Tab ``Catalogo'' nell'area cliente.

\textbf{Perché esiste:}
Consente al cliente di visionare offerta senza rischio di modifica accidentale.

\subsection{OperazioniClientePanel -- File: OperazioniClientePanel.java (Cliente)}

\textbf{Cosa fa:}
Dual-table view: storico vendite (quando, quale veicolo, prezzo) e storico test drive (quando, quale veicolo, eliminabili).
Cliente può cancellare prenotazioni test drive future.

\textbf{Come viene usato:}
Tab ``Le Mie Operazioni'' nell'area cliente.

\textbf{Perché esiste:}
Consente al cliente di tracciare la propria cronologia e gestire prenotazioni.

\section{Livello Control: Orchestrazione della Logica}

\subsection{ConcessionarioController() -- Constructor. File: ConcessionarioController.java}

\textbf{Cosa fa:}
Istanzia tutti i DAO impl (ClienteDAOImpl, VeicoloDAOImpl, VenditaDAOImpl, TestDriveDAOImpl, DipendenteDAOImpl).

\textbf{Come viene usato:}
Constructor chiamato dalla GUI (ConcessionarioFrame, pannelli).

\textbf{Perché esiste:}
Il controller è il depositario di tutti i DAO. Da qui, GUI accede ai dati senza sapere quale DAO concreto sta usando (polimorfismo).

\subsection{ConcessionarioController.autenticaCliente() -- File: ConcessionarioController.java}

\textbf{Cosa fa:}
Riceve email e password da GUI e ritorna un oggetto Cliente se autenticazione ha successo.

Logica interna:
\begin{enumerate}
  \item Controlla che email non sia null e non sia stringa vuota. Se fallisce, solleva \texttt{ValidationException}.
  \item Controlla che password non sia null e non sia stringa vuota. Se fallisce, solleva \texttt{ValidationException}.
  \item Valida formato email via \texttt{InputValidator.isValidEmail(email)}: regex controlla che ci sia '@' e '.' nel posto giusto.
    Se email non ha formato valido, solleva \texttt{ValidationException(``Email non valida'')}.
  \item Chiama \texttt{ClienteDAO.autenticaCliente(email, password)}.
    Il DAO fa una query case-insensitive (mario@example.com = MARIO@EXAMPLE.COM) e confronta password.
  \item Se DAO ritorna Cliente non-null, lo ritorna al caller (GUI). Autenticazione riuscita.
  \item Se DAO ritorna null, solleva \texttt{ValidationException(``Email o password non corretti'')}.
\end{enumerate}

\textbf{Dettagli di implementazione interna:}

La validazione è doppia:
\begin{itemize}
  \item \textbf{Nel Controller:} formato email, non-null check.
  \item \textbf{Nel DAO:} query al DB con case-insensitive, confronto password.
\end{itemize}

Perché separare? Se la GUI è buggy e passa email malformata, il Controller la butta fuori senza neanche toccare il DB (economia).
Se invece API REST chiama il Controller, la validazione avviene lo stesso (non dipendi dalla GUI).

\textbf{Edge case: email diverso da maiuscole/minuscole}

Scenario: utente registrato come \texttt{Mario@Example.com}, prova login con \texttt{mario@example.com}.
Il DAO usa \texttt{LOWER(email)} in query SQL, quindi entrambe matchano. Bene.

Scenario: utente registrato come \texttt{Mario@Example.com}, prova login con \texttt{mario@EXAMPLE.COM}.
Idem: LOWER() gestisce entrambe. Bene.

\textbf{Come viene usato:}
GUI (handleClienteLogin in ConcessionarioFrame) estrae email e password da form, chiama questo metodo.
Se succeed, ritorna Cliente con id/nome/email. La GUI aggiorna header (``Autenticato come Mario'').
Se exception, GUI catcha e mostra dialog.

\textbf{Perché esiste:}
Centralizza autenticazione. Se in futuro cambi il meccanismo (da plaintext password a bcrypt-hashed), cambi SOLO qui e nel DAO.
La GUI non sa che cosa è successo: chiama il controller, riceve Cliente o exception, basta.

\subsection{ConcessionarioController.autenticaDipendente() -- File: ConcessionarioController.java}

\textbf{Cosa fa:}
Analogo a autenticaCliente, ma riceve ID (intero) e password.
Valida non-null, chiama \texttt{DipendenteDAO.autenticaDipendente()}.

\textbf{Come viene usato:}
handleDipendenteLogin nel frame.

\subsection{ConcessionarioController.inserisciCliente() -- File: ConcessionarioController.java}

\textbf{Cosa fa:}
Inserisce un nuovo cliente nel DB. Riceve nome, tipo, email, telefono, password.

Logica interna (validazione a strati):

\begin{enumerate}
  \item \textbf{Validazione nome:} null o empty $\Rightarrow$ \texttt{ValidationException}.
  \item \textbf{Validazione email:} null, empty, o formato invalido $\Rightarrow$ \texttt{ValidationException}.
    Es. \texttt{mario} senza '@' fallisce. \texttt{mario@} fallisce (no nome dominio). \texttt{mario@domain} fallisce (no TLD).
\subsection{InputValidator.isValidPhone()} -- File: InputValidator.java}

\textbf{Cosa fa:}
Valida numero di telefono: esattamente 10 cifre numeriche.
Regex: \texttt{\^\\d\{10\}\$}.

\textbf{Regex Breakdown:}

\begin{verbatim}
^      # Inizio stringa
\\d{10} # Esattamente 10 cifre (0-9)
\$      # Fine stringa
\end{verbatim}

\texttt{\\d} = qualunque cifra da 0 a 9.
\texttt{\\d\{10\}} = ripeti esattamente 10 volte.

\textbf{Esempi che PASSANO:}

\begin{itemize}
  \item \texttt{0123456789}: 10 cifre. ✓
  \item \texttt{1234567890}: 10 cifre. ✓
  \item \texttt{9999999999}: 10 noni. ✓
\end{itemize}

\textbf{Esempi che FALLANO:}

\begin{itemize}
  \item \texttt{123456789}: 9 cifre, < 10. ✗
  \item \texttt{12345678901}: 11 cifre, > 10. ✗
  \item \texttt{123-456-7890}: Contiene trattini. Il trattino non è \\d. ✗
  \item \texttt{123 456 7890}: Contiene spazi. Lo spazio non è \\d. ✗
  \item \texttt{+39 123456789}: Contiene '+' e spazio. ✗
\end{itemize}

\textbf{Nota sul Design:}

Questo regex è rigido: accetta SOLO 10 cifre consecutive, niente altro.

Reale: numeri telefonici possono avere formati diversi (con trattini, spazi, +codice paese).

Per questo progetto, "10 cifre" è sufficiente ed è facile da validare.

\textbf{Come viene usato:}

Controller.inserisciCliente() valida campo telefono prima di passare al DAO.

\subsection{InputValidator.isValidPlate()} -- File: InputValidator.java}

\textbf{Cosa fa:}
Controlla targa: 5-10 caratteri alfanumerici.

\textbf{Come viene usato:}
VeicoliPanel.onAggiungi() valida targa.

\subsection{InputValidator.normalizePlate()} -- File: InputValidator.java}

\textbf{Cosa fa:}
Trim e UPPER: \texttt{``ab123cd''} diventa \texttt{``AB123CD''}.

\textbf{Come viene usato:}
VeicoliPanel: prima di INSERT o search, normalizza targa per coerenza.

\subsection{InputValidator.parseDateOrThrow()} -- File: InputValidator.java}

\textbf{Cosa fa:}
Parsing stringa a LocalDate in formato ISO (YYYY-MM-DD).
Se fallisce, solleva exception.

\textbf{Come viene usato:}
TestDrive form picker: data inserita dall'utente, convertita a LocalDate.

\subsection{InputValidator.requireTodayOrFuture()} -- File: InputValidator.java}

\textbf{Cosa fa:}
Controlla che data non sia nel passato.

\textbf{Come viene usato:}
Controller.richiediTestDrive() valida che data è oggi o futura.

\subsection{UiSupport.wrapTitled()}} -- File: UiSupport.java}

\textbf{Cosa fa:}
Riceve un JPanel, lo confeziona in un wrapper con titolo borderizzato.

\textbf{Come viene usato:}
Pannelli l'usano per marcare zone del form (``Dettagli Cliente'', ``Indirizzo'', etc.).

\subsection{UiSupport.showErr()}} -- File: UiSupport.java}

\textbf{Cosa fa:}
Mostra un JOptionPane. dialog Error con messaggio.

\textbf{Come viene usato:}
Ovunque in GUI, quando catchi una exception.

\subsection{UiSupport.toUserMessage()}} -- File: UiSupport.java}

\textbf{Cosa fa:}
Riceve un'exception (ValidationException, DatabaseException), ritorna un messaggio amichevole.

Traduce messaggi tecnici:
\begin{itemize}
  \item ``cliente\_email\_key'' (constraint violation) $\to$ ``Email già registrata: usa un'altra''.
  \item ``limit\_testdrive'' $\to$ ``Hai raggiunto il limite di 3 test drive al mese''.
  \item Fallback: ritorna exception.getMessage().
\end{itemize}

\textbf{Come viene usato:}
\texttt{try \{ ... \} catch (ValidationException|DatabaseException ex) \{ UiSupport.showErr(this, UiSupport.toUserMessage(ex)); \}}.

\section{Flusso Completo di un'Operazione Tipica}

Esempio: Dipendente aggiunge un nuovo cliente.

\begin{enumerate}
  \item \textbf{Utente clicca bottone ``Aggiungi'' in ClientiPanel.}
  \item \textbf{ClientiPanel.onAggiungi() esegue:}
  \begin{enumerate}
    \item Legge campi form (nome, email, telefono, tipo).
    \item Chiama \texttt{Controller.inserisciCliente(nome, tipo, email, telefono, password)}.
  \end{enumerate}
  
  \item \textbf{Controller.inserisciCliente() esegue:}
  \begin{enumerate}
    \item Valida nome non vuoto.
    \item Valida email con \texttt{InputValidator.isValidEmail(email)}.
    \item Valida telefono con \texttt{InputValidator.isValidPhone(telefono)}.
    \item Controlla duplicato: \texttt{ClienteDAO.getClienteByEmail(email)}.
    \item Se trovato, solleva \texttt{ValidationException(``Email già registrata'')}.
    \item Se ok, chiama \texttt{ClienteDAO.inserisciCliente(nome, tipo, email, telefono, password)}.
  \end{enumerate}
  
  \item \textbf{ClienteDAOImpl.inserisciCliente() esegue:}
  \begin{enumerate}
    \item Apre connessione JDBC via \texttt{DatabaseManager.getInstance().getConnection()}.
    \item Genera SQL: \texttt{INSERT INTO Cliente (nome, tipo, email, telefono, password) VALUES (?, ?, ?, ?, ?)}.
    \item Imposta parametri con \texttt{ps.setString()}.
    \item Esegue \texttt{ps.executeUpdate()}.
    \item Se eccezione SQL (email duplicata, constraint error), wrappa in \texttt{DatabaseException}.
    \item Chiude connessione (try-with-resources).
  \end{enumerate}
  
  \item \textbf{Database PostgreSQL:}
  \begin{enumerate}
    \item Riceve INSERT.
    \item Controlla constraint: email UNIQUE, telefono formato, etc.
    \item Se ok, inserisce riga, ritorna success.
    \item Se email non unique, riga già esiste $\Rightarrow$ exception.
  \end{enumerate}
  
  \item \textbf{Ritorno alla GUI:}
  \begin{enumerate}
    \item Se success: ClientiPanel.onAggiungi() mostra dialog ``Successo'', chiama \texttt{loadClienti()} per ricaricare tabella, clear form.
    \item Se exception: catch \texttt{ValidationException|DatabaseException ex}, traduce con \texttt{UiSupport.toUserMessage(ex)}, mostra error dialog.
  \end{enumerate}
  
  \item \textbf{Utente vede:}
  \begin{enumerate}
    \item Successo: ``Cliente aggiunto con successo'', tabella aggiornata con nuovo cliente.
    \item Errore: ``Email già registrata: usa un'altra''.
  \end{enumerate}
\end{enumerate}

\section{Principi Architetturali Applicati}

\subsection{Separation of Concerns}

Ogni livello ha una responsabilità:
\begin{itemize}
  \item Boundary (GUI): widget, layout, event handling.
  \item Control (Controller): validazione, logica business, coordinamento.
  \item DAO: SQL, JDBC, eccezioni DB.
  \item Entity: dati pure (immutabili).
  \item DB: integrità, trigger, funzioni computate.
\end{itemize}

Cambio nel DB (schema, funzione) non tocca la GUI.
Cambio nella GUI (layout) non tocca la logica.

\subsection{Polymorphism via Interfacce}

Controller dipende da \texttt{ClienteDAO} (interfaccia), non da \texttt{ClienteDAOImpl}.
Questo consente:
\begin{itemize}
  \item Swap implementazione (PostgreSQL $\to$ MongoDB).
  \item Test con MockClienteDAO (simula DB senza connessione reale).
\end{itemize}

\subsection{Immutability Entities}

Cliente, Dipendente, Veicolo hanno campi final.
Vantaggio: thread-safe per natura. Non rischio race condition di stato entity.

\subsection{Try-with-Resources}

Tutte le connessioni JDBC aperte con try-with-resources.
Garantisce: anche se exception, risorsa è chiusa. Niente memory leak.

\subsection{Exception Hierarchy}

Eccezioni specializzate:
\begin{itemize}
  \item \texttt{ConcessionarioException}: base.
  \item \texttt{ValidationException}: estende, per errori input.
  \item \texttt{DatabaseException}: estende, per errori DB.
\end{itemize}

GUI catcha in diversi modi: \texttt{catch (ValidationException ex)} vs. \texttt{catch (DatabaseException ex)}.

\textbf{Ultimo consiglio:} Leggi questo documento e il codice in parallelo. Quando vedi una classe nel codice, guarda qui cosa fa. Chiedi sempre: ``Perché esiste questa funzione?'' Una buona architettura risponde questa domanda chiaramente.

\end{document}
